% ─────────────────────────────────────────────────────────────────────────────
% Section I — Introduction
% Paper: Scaling Laws for Cyber Resilience: Modeling the Persistence of
%        Backdoors in Post-trained Large Language Models
% ─────────────────────────────────────────────────────────────────────────────

\section{Introduction}
\label{sec:intro}

Large language models (LLMs) have demonstrated remarkable capabilities across
security-sensitive domains, including penetration testing assistance~\cite{happe2023got},
vulnerability analysis~\cite{pearce2023examining}, and autonomous cyberdefense~\cite{kang2024exploiting}.
This proliferation, however, introduces a critical attack surface: adversaries
may embed \emph{backdoors}---hidden behavioral triggers that cause targeted
model outputs when activated---during the supply chain of model training.
Unlike classical software vulnerabilities, LLM backdoors survive post-training
alignment procedures (supervised fine-tuning, SFT, and reinforcement learning from
human feedback, RLHF), creating a persistent threat even after safety hardening.

A fundamental open question concerns the relationship between \emph{model scale}
and \emph{backdoor persistence}: do larger models resist fine-tuning-based
backdoor removal more than smaller ones? Empirical evidence from prior
work~\cite{yang2024watch,chen2021badnl,wang2024rlhf_backdoor} suggests a positive
correlation between parameter count and attack success rate (ASR) after RLHF,
but no formal metric captures this relationship in a principled way suitable
for risk quantification or regulatory compliance.

\textbf{Contributions.} This paper makes four concrete contributions:

\begin{enumerate}
  \item We propose the \emph{Cyber Resilience Scaling Coefficient} (CRSC), a
  closed-form metric that quantifies backdoor persistence after fine-tuning as a
  function of model scale $N$, poison rate $\rho$, and fine-tuning dataset $D_\text{ft}$
  (Section~\ref{sec:threat}).

  \item We prove two theoretical properties: (i) monotonicity of CRSC in $N$
  for fixed $(\rho, D_\text{ft})$ --- larger models are provably harder to
  sanitize (Proposition~\ref{prop:monotone}); and (ii) CRSC is a concentration-of-measure
  proxy for ASR with bounded estimation error (Proposition~\ref{prop:proxy}).

  \item We design a three-agent evaluation pipeline (Detection, Analysis, Response)
  and validate CRSC on synthetic data calibrated against the TrojAI NIST 2024
  benchmark~\cite{trojai2024}, achieving F1 $= 0.891$ and AUC-ROC $= 0.924$ at
  $\tau = 0.62$ (Section~\ref{sec:experiments}).

  \item We map CRSC to the MITRE ATLAS adversarial ML framework and the
  NIST AI Risk Management Framework (AI RMF), providing actionable mitigation
  guidance for practitioners (Section~\ref{sec:security}).
\end{enumerate}

\textbf{Scope and limitations.} Our analysis assumes white-box access to
hidden states (a common red-team setting~\cite{perez2022red}), static trigger
patterns, and transformer architectures. Adaptive triggers and black-box
settings remain open challenges; we discuss them in Section~\ref{sec:discussion}.

\textbf{Organization.} Sections~\ref{sec:background}--\ref{sec:security} cover
background, threat model, pipeline, and compliance. Section~\ref{sec:experiments}
reports experimental results; Sections~\ref{sec:discussion}--\ref{sec:conclusion}
discuss implications and conclude.
